\section{Why the Standard Model: Density, Ramanujan, Probability}\label{sec:why_sm}

The landscape analysis (Section~\ref{sec:landscape}) established that the Standard Model's non-simple product structure is the generic outcome of the variational principle. But \textbf{why SU(2) $\times$ SU(3) specifically}? The answer lies in the intersection of three independent criteria: density, spectral optimality (Ramanujan property), and accessibility.

\subsection{The Three-Criteria Argument}\label{sec:three_criteria}

The Standard Model is selected by the confluence of three independent filters:

\begin{enumerate}
\item \textbf{Killing + $T_{\mathrm{cl}}$} $\to$ favors CLOSED\_NON\_SIMPLE (53.1\%)
\item \textbf{Ramanujan} $\to$ selects among CLOSED\_NON\_SIMPLE the spectrally optimal
\item \textbf{Accessibility} $\to$ SU(2) $\times$ SU(3) is the CLOSED\_NON\_SIMPLE most accessible
\end{enumerate}

\subsubsection{Criterion 1: Topological Dominance}

The landscape measure shows that non-simple products dominate (53.1\% of seeds). This is the first filter: the universe must select among non-simple structures.

\subsubsection{Criterion 2: The Ramanujan Filter}

Among all CLOSED\_NON\_SIMPLE structures, which is spectrally optimal? The Ramanujan property---optimal spectral expansion of the Cartan--Weyl root graph---provides the answer.

\begin{theorem}[Ramanujan Property of SM Factors]
The CW graphs of $\su(2)$ and $\su(3)$ are Ramanujan with ratios $\mathcal{R}(\su(2)) = 1/2$ and $\mathcal{R}(\su(3)) = (\sqrt{57} - 3)/(2\sqrt{17}) \approx 0.5523$, respectively. Both pass all six alternative definitions of the Ramanujan property unanimously.
\end{theorem}

The product $\su(3) \oplus \su(2)$ is also Ramanujan (ratio 0.657), confirming P7: Ramanujan $\times$ Ramanujan $\to$ Ramanujan.

F$_4$ is Ramanujan but does not emerge blindly ($P < 3.1\%$). The SM combination is \textbf{Ramanujan + accessible}---the only one satisfying both criteria.

\subsubsection{Criterion 3: Accessibility}

Among all CLOSED\_NON\_SIMPLE structures, SU(2) $\times$ SU(3) is the most accessible: it has the largest basin of attraction in the extended phase space, and it is the unique product of simple algebras that is both Ramanujan and has a large enough basin to be found from random initialization.

\subsection{The Density Argument}\label{sec:density_argument}

The density phase transition (Section~\ref{sec:density}) provides the mechanistic bridge between the variational principle and the SM selection.

For SU(4) with 15 generators, the step function at $n_{\mathrm{gen}} = 15$ means that the SM gauge group emerges \textbf{only} when the relational graph reaches the critical density $\rho_c = 1$. Below this threshold: no algebraic structure. At this threshold: exact SU(4) emerges.

The SM gauge group SU(3) $\times$ SU(2) $\times$ U(1) has dimension $8 + 3 + 1 = 12$. Within the SU(4) generator budget (15 generators), the SM fits comfortably with $12 \leq 15$, leaving 3 generators for the U(1) hypercharge direction and the trace component.

\subsubsection{The SM Lives in the High-Density Region}

The SM is not a marginal case. It lives in the region of \textbf{high relational density}---well above the percolation threshold. This is not accidental: the SM gauge group has exactly the right dimension to emerge at the density where the universe's relational graph crystallized.

\subsection{The Multi-Layer Probability Argument}\label{sec:multi_layer}

A six-layer convergent argument establishes not only that the SM-type structure is generic, but that SU(3) $\times$ SU(2) $\times$ U(1) is the \textbf{unique} spectrally optimal choice:

\begin{equation}
\begin{array}{lcc}
\hline
\textbf{Layer} & \textbf{Mechanism} & \textbf{Value} \\
\hline
1 & \text{Topological dominance: } P(\text{non-simple}) & 53.1\% \\
2 & \text{Discrete attractor: dominant band B}_2 & 53.4\% \\
3 & \text{Ramanujan filter: all dim $\leq 51$ ORDERED} & 100\% \\
4 & \text{Spectral uniqueness: SM = argmin($\Sigma$R)} & \text{\#1/100 pairs} \\
5 & \text{D1 verification: SM product is 6/6 Ramanujan} & 6/6 \\
6 & \text{Composability P7: Ramanujan $\times$ Ramanujan $\to$ Ramanujan} & \text{Unanimous} \\
\hline
\end{array}
\end{equation}

The physical constraints are: (a) at least one rank $\geq 2$ factor (color), (b) at least one rank 1 non-abelian factor (weak isospin), (c) U(1) hypercharge (trace constraint). Under these constraints, $\su(3) \oplus \su(2)$ achieves the \textbf{unique minimum} of the total Ramanujan cost $\Sigma R(h_i)$ among all 100 viable color $\times$ weak pairs.

\subsection{Conclusion: Why the SM}

The Standard Model gauge group is not an arbitrary choice. It is the \textbf{unique} outcome of three independent filters acting in concert:

\begin{enumerate}
\item The variational principle favors non-simple products (53.1\%)
\item Among these, the Ramanujan property selects SU(3) $\times$ SU(2) as spectrally optimal
\item Among the Ramanujan candidates, SU(3) $\times$ SU(2) has the largest basin of attraction
\end{enumerate}

The SM is the \textbf{only} gauge group that satisfies all three criteria simultaneously. No other combination comes close: the next-best option has $> 241\%$ error in dual Coxeter numbers.
